\documentclass[11pt,a4paper]{article}
\usepackage[T1]{fontenc}
\usepackage[utf8]{inputenc}
\usepackage{amsmath,amssymb,xcolor}
\usepackage[margin=25mm]{geometry}
\usepackage[colorlinks=true,urlcolor=blue]{hyperref}
\definecolor{nexorablue}{HTML}{2459B5}
\setlength{\parindent}{0pt}
\setlength{\parskip}{7pt}
\setlength{\emergencystretch}{2em}
\renewcommand{\familydefault}{\sfdefault}
\begin{document}

{\small\bfseries\color{nexorablue} NEXORA CPGE}\par
{\small MPSI / NEX-MES-01 / Énoncé}\par
{\LARGE\bfseries Quatre lectures au banc optique}\par
\vspace{0.25em}\hrule\vspace{0.6em}

Quatre lectures indépendantes d’une position donnent $20{,}14$, $20{,}16$, $20{,}17$ et $20{,}13\,\mathrm{cm}$. Calculer la moyenne, l’écart-type expérimental et l’incertitude-type de la moyenne en ne tenant compte que des répétitions.


\end{document}
